% Immutable cited-input supplement. Target proofs remain embargoed.
% Target: night:P4.21; tag 02WT.
% Frozen source SHA256 74497bfe327fae25c9cdcb17334b3bf42a9793a9df7c0a4404d97568f2f760f0.
% CompendiumRef is resolved against the cited S6_09 source edition, not silently identified with the current series statement.

% Exact cited input: series:II / night:P4.12
\begin{proposition}[P4.12]\hypertarget{tag-02WJ}{}\label{night:P4.12}\label{stmt:P4.12}


In the standing torus-fibration setting of this section, let $f^{-1}(b)=mS$
have irreducible reduced support. Then
\[
 \omega_S^{\otimes m}\simeq\OO_S,\qquad
 \chi(\OO_S)=0,\qquad h^1(\OO_S)\in\{1,2\}.
\]
Every normal such support is smooth. For a singular, necessarily
nonnormal support, a minimal resolved model of its normalization
has Kodaira dimension $-\infty$.
\end{proposition}

% Exact cited input: series:II / night:P4.13
\begin{proposition}[P4.13]\hypertarget{tag-02WK}{}\label{night:P4.13}\label{stmt:P4.13}


Let a proper holomorphic degeneration of two-dimensional complex
tori have smooth total threefold and central fibre $mS$ with
irreducible reduced support. If $A$ is the monodromy on $H_1$ of
a nearby torus, then
\[
 m e(S_{\mathrm{reg}})=\det(I-A^m),\qquad
 0\le e(S_{\mathrm{reg}})\le16/m.
\]
This assertion requires no Betti-number hypothesis on a compactification.
\end{proposition}

% Exact cited input: series:II / night:P4.20
\begin{proposition}[P4.20]\hypertarget{tag-02WS}{}\label{night:P4.20}\label{stmt:P4.20}


Let $mS$ be an irreducible fibre of a proper holomorphic
degeneration of complex two-tori with smooth total threefold.
Suppose the normalization of $S$ is smooth and, outside
finitely many points, $S$ is smooth or has the ordinary-double
germ $xy=0$. Then $S$ has normal crossings everywhere.
In the global setting of 4.2, P4.6--P4.8 apply to this subclass.
\end{proposition}

% Exact cited input: series:II / night:P4.6
\begin{proposition}[P4.6]\hypertarget{tag-02XY}{}\label{night:P4.6}\label{stmt:P4.6}


Let $X$ be a compact complex threefold with $b_2(X)=0$, and let $F$ be a
reduced compact normal-crossings divisor. We allow different local
branches of one global irreducible component to meet. If $T(F)$ is the
number of ordinary triple points, then
\[
 e(F)=T(F),\qquad T(F)\text{ is even}.
\]
This assertion concerns the reduced support, regardless of the
multiplicities with which it occurs in a fibre.
\end{proposition}

% Exact cited input: series:II / night:P4.8
\begin{proposition}[P4.8]\hypertarget{tag-02Y0}{}\label{night:P4.8}\label{stmt:P4.8}


For a torus fibration in P4.7 whose reduced supports have normal
crossings, an irreducible singular support with
rational normalization has Euler number two; every other irreducible
support has Euler number zero. In the zero case the supports have the
forms stated in the prompt: a smooth support has a finite unramified
torus cover; a singular irreducible support has smooth minimal
elliptic-ruled or Hopf normalization; a reducible support has smooth
minimal Hopf or non-Kähler elliptic normalization components, with
smooth elliptic double curves and finite unramified branch maps.
Thus $c_3[X]=0$ precisely when every fibre belongs to these zero cases.
\end{proposition}

% Exact cited input: series:II / night:P4R.1
\begin{proposition}[P4R.1]\hypertarget{tag-02Y2}{}\label{night:P4R.1}\label{stmt:P4R.1}


For a proper holomorphic degeneration of complex two-tori with
smooth total threefold and central fibre $mS$, where $S$ is
irreducible and reduced, the group $\pi_1(S)$ has an abelian
normal subgroup with cyclic quotient of order dividing $m$.
\end{proposition}

% Exact cited input: series:II / night:P4R.3
\begin{proposition}[P4R.3]\hypertarget{tag-02YF}{}\label{night:P4R.3}\label{stmt:P4R.3}


In the standing torus-fibration setting of this section, let an irreducible singular fibre
support $S$ have smooth normalization $\nu:T\to S$ whose
minimal model is ruled over a curve of genus $g\ge2$.
Write $D$ for the conductor divisor on $T$ and
$C=\operatorname{Sing}(S)_{\rm red}$.
The inclusion $D_{\rm red}\to T$ induces an isomorphism on
fundamental groups, and the normalization pushout induces
$\pi_1(S)\simeq\pi_1(C)$. Moreover,
\[
 e(C)\ge0,\qquad e(S)>0.
\]
The assertion includes conductor components on different
ruling fibres that map onto the same curve in $C$.
\end{proposition}

